\chapter{L'evoluzione dell'argomentazione critica}
\label{cap:evoluzione-argomentazione}

\todoscrivi{Nucleo della tesi (ADR-0007, 2026-09-14). Rassegna critica organizzata
CRONOLOGICAMENTE, non per temi: e' la forma esplicitamente richiesta dal relatore
("come si e' evoluta nel tempo"). Ogni sezione e' una fase storica; dentro ciascuna
puo' comparire piu' di un filone tematico se sono emersi nello stesso periodo. Le
schede complete dei paper vivono in docs/literature/; qui entra solo cio' che serve
a costruire l'argomento. Fonti gia' individuate (da leggere per intero prima di
citarle) in docs/literature/da-recuperare.md, quattro filoni dell'11 agosto 2026:
estetica/casualita', locus della creativita', falsa rivendicazione/etica,
training-vs-ispirazione.}

\section{Metodologia della rassegna}
\label{sec:metodologia-revisione}
% Database interrogati, stringhe di ricerca, criteri di inclusione ed esclusione,
% intervallo temporale, numero di lavori esaminati e selezionati. Rende la
% rassegna riproducibile e la difende in sede di discussione.

\section{2017: la rivendicazione esplicita e la prima obiezione estetica}
\label{sec:evo-2017}
% Punto di partenza: Elgammal et al. (2017), le Creative Adversarial Networks
% (gia' presentate come caso in \cref{sec:can}). Prima risposta teorica, filone
% "estetica/casualita'": l'estetica simbolica di Susanne Langer (l'arte come
% forma imbevuta di emozione) e la letteratura correlata (Qi 2019, Cassirer se
% serve la genealogia del concetto). Nodo da trattare con precisione (gia'
% presente in thesis/extra/proposta-relatore.tex): la critica NON e' "e' solo
% casualita'" -- la creativita' computazionale si definisce proprio in
% opposizione al caso, liquidarla cosi' sarebbe un errore facile da rilevare.

\section{Il locus della creativita': algoritmo o uso umano?}
\label{sec:evo-locus}
% Filone "locus della creativita'": il dibattito si sposta da "l'algoritmo e'
% creativo?" a "chi/cosa e' creativo nel sistema uomo+AI?". Letteratura sulla
% co-creativita' uomo-AI (2024) e sull'identita' creativa di chi usa lo
% strumento. Collegamento a D-024 (l'accessibilita' come beneficio indipendente
% dalla legittimita' del training) senza fondere i due argomenti.

\section{L'ondata etica: antropomorfizzazione e le sue conseguenze}
\label{sec:evo-etica}
% Filone "falsa rivendicazione come problema etico": con l'esplosione dei
% modelli diffusivi (DALL-E, Midjourney, Stable Diffusion, Nano Banana) il
% dibattito diventa mainstream. Attribuire creativita' a un sistema che non la
% possiede e' antropomorfizzazione -- non un vezzo linguistico, distorce il
% giudizio morale (valutazione del pubblico, prezzo/percezione del lavoro degli
% artisti umani, mercato, regolazione). Include l'obiezione training-vs-
% ispirazione (il problema del symbol grounding, il dibattito sul copyright: se
% manca la componente emotiva nell'output, manca anche nel training, e
% l'analogia con l'ispirazione umana perde fondamento -- gia' argomentato in
% thesis/extra/proposta-relatore.tex).

\section{2025-2026: lo stato piu' recente del dibattito}
\label{sec:evo-recente}
% Le fonti piu' aggiornate reperibili al momento della stesura: rassegne sulla
% creativita' computazionale applicata al machine learning (Franceschelli \&
% Musolesi, ACM Computing Surveys 2024), critiche alle metriche di valutazione
% della creativita' esistenti, misure quantitative di creativita' per modelli
% generativi visivi, lavori sull'agentivita' intenzionale (\emph{Creativity
% Reconsidered}, 2026). Da verificare quali, lette per intero, criticano
% esplicitamente la rivendicazione delle CAN o la confermano.

\section{Lacune e domanda di ricerca}
\label{sec:gap}
% La sezione piu' importante del capitolo: da qui deve discendere la domanda di
% ricerca dichiarata in \cref{sec:domande-ricerca}. Il lavoro DCGAN/CAN/E8 di
% Gian NON entra qui come prova (ADR-0007): puo' essere nominato al piu' nelle
% conclusioni come sviluppo futuro/parallelo.
